\chapter{The Cohomological Type Stack}

A mature foundational project feels unified when its central ideas are not introduced as
separate tricks for separate subproblems, but as a small vocabulary that keeps reappearing
until the reader realizes that the whole enterprise is one object seen under several
illuminations. That is one reason some mathematical projects feel inevitable rather than
accidental. They persuade the reader that once a few basic constructions are granted, the
later chapters are not bolt-ons but consequences. \DepPy{} should be read in that spirit.

The repository is easiest to trivialize when its components are described in isolation:
there is a plain type system; there is a hybrid extension; there are trust levels; there is
an oracle monad; there is a cohomological bug story; there are theory packs; there is a
zero-change verifier; there is a generation agent. Read that way, the repository looks like
a mixed research notebook. Read more carefully, it looks like a single stack whose levels
share one grammar. The grammar is the grammar of local sections, transport, gluing,
obstruction, repair, and eventual assembly.

\section*{The one-stack principle}
The most useful way to say this is the \emph{one-stack principle}:
\begin{quote}
\emph{Every major subsystem in \DepPy{} is a manifestation of one cohomological type stack.
The lower levels supply the objects and local obligations; the middle levels supply the
witnesses, trust annotations, and proof boundaries; the upper levels supply the methods by
which typed local material is searched, transported, rendered, and assembled.}
\end{quote}

The stack is not a metaphor pasted over arbitrary code. It is a practical account of how
the repository is organized.

\begin{definition}[The cohomological type stack]
The cohomological type stack is the five-level organization
\[
  \mathsf{Foundation} \to \mathsf{Object} \to \mathsf{Evidence} \to \mathsf{Semantics} \to \mathsf{Assembly}
\]
in which:
\begin{enumerate}[leftmargin=2em, itemsep=0.35em]
  \item \textbf{Foundation} supplies sites, covers, restriction maps, and presheaf structure.
  \item \textbf{Object} supplies the exact type language of carriers, refinements, and
  dependent formers.
  \item \textbf{Evidence} lifts those objects into hybrid witnesses distributed across
  intent, structure, semantics, proof, and code.
  \item \textbf{Semantics} classifies and composes evidence via trust lattices, oracle
  effects, theorem translation, and residual obligations.
  \item \textbf{Assembly} grows, repairs, and glues local artifacts into runnable or
  checkable systems.
\end{enumerate}
\end{definition}

The reward for naming the stack explicitly is that vocabulary starts to stabilize. A site is
a place where one may speak locally. A section is typed local evidence. Transport is the
movement of that evidence along a semantic map. Gluing is the passage from local evidence to
a global artifact. Obstruction is non-gluing. Repair is a refinement of the cover, the local
claim, or the evidence attached to the claim. The same words remain meaningful whether one is
speaking about subtyping, hallucination, theory-pack discharge, or multi-module synthesis.

\section*{Five manifestations of one vocabulary}
The repository's apparent diversity can now be restated as a controlled repetition of this
vocabulary.

At the foundational level, one has sites and presheaves. A function is not seen as one flat
blob of semantics but as a cover of boundaries, guards, intermediate values, call sites, and
error-sensitive observations. That is already a local-to-global stance.

At the object level, one has the plain language in \repo{src/deppy/types/}. Here the
repository defines the actual symbolic objects that higher levels will manipulate:
\py{TypeBase}, refinement types, $\Pi$-types, $\Sigma$-types, equality-like constructions,
indexed families, and subtype obligations. This level says what must be true.

At the evidence level, one has the hybrid language in \repo{src/deppy/hybrid/core/}. Here
those same obligations acquire multiple manifestations: intent text, structural predicates,
semantic rubrics, proof artifacts, and code. This level says how a claim is witnessed and
how much confidence should be placed in the witness.

At the semantic level, one has trust, oracle effects, proof translation, discharge, and the
bookkeeping of residual honesty. This level says what kind of evidence one has, what it is
worth, and how it composes.

At the assembly level, one has type expansion, theory packs, zero-change checking, and the
generation agent. This level says how local pieces are found, prioritized, and glued into a
coherent system.

The central argument of this chapter is that these are not five separate research projects.
They are five manifestations of one project.

\begin{figure}[h]
\centering
\begin{tikzpicture}[
  every node/.style={align=center},
  box/.style={draw=chapterblue, rounded corners=2pt, fill=chapterblue!6, minimum width=10.6cm, minimum height=0.95cm},
  note/.style={draw=accentgold, rounded corners=2pt, fill=accentgold!12, minimum width=3.0cm, minimum height=0.8cm},
  >=Latex
]
\node[box] (foundation) {Foundation \quad sites, covers, restrictions, presheaves};
\node[box, below=0.5cm of foundation] (object) {Object \quad \repo{src/deppy/types/}: carriers, refinements, dependent formers, subtype obligations};
\node[box, below=0.5cm of object] (evidence) {Evidence \quad \repo{src/deppy/hybrid/core/}: layered witnesses across intent, structure, semantics, proof, and code};
\node[box, below=0.5cm of evidence] (semantics) {Semantics \quad trust lattices, oracle monad, Lean translation, discharge, residuals};
\node[box, below=0.5cm of semantics] (assembly) {Assembly \quad theory packs, zero-change analysis, ideation, type expansion, agent synthesis};
\node[note, right=1.2cm of evidence] (glue) {same vocabulary \ section / transport / glue / obstruct / repair};
\foreach \a/\b in {foundation/object,object/evidence,evidence/semantics,semantics/assembly} {
  \draw[->, thick, chapterblue] (\a) -- (\b);
}
\draw[->, thick, accentgold] (object.east) -- (glue.west);
\draw[->, thick, accentgold] (assembly.east) .. controls +(1.2,0.1) and +(0.0,-0.9) .. (glue.south);
\end{tikzpicture}
\caption{The cohomological type stack. The project becomes conceptually unified once the same local-to-global vocabulary is seen to govern each level.}
\end{figure}

\section*{Why the plain and hybrid systems belong to the same project}
The plain and hybrid systems are sometimes described as though they belonged to different
ages of the repository. This is historically understandable and mathematically misleading.
The plain system is not an abandoned first draft; it is the exact object language that makes
the hybrid system possible. The hybrid system is not a repudiation of the plain system; it
is the evidence-bearing lift of the same obligations into a larger universe.

That relation can be stated succinctly. The plain system determines what a value, function,
refinement, or witness \emph{is}. The hybrid system determines what it means to possess
multiple, potentially nonidentical, forms of evidence about that object and to compare those
forms without erasing their differences. The hybrid system therefore presupposes the plain
one in much the way a semantics of paths presupposes points.

\section*{One lexicon across engineering concerns}
Once the stack is visible, many engineering decisions in the repository become easier to
justify.

The mixed-mode syntax is no longer an ergonomic front-end awkwardly bolted onto a verifier.
It becomes the public boundary at which ordinary Python enters the stack. Surface forms such
as \py{@guarantee}, \py{assume()}, \py{check()}, and \py{hole()} are user-facing names for
moves that later become sections, obligations, and proof targets.

Zero-change verification is no longer merely a convenience mode. It becomes the inverse
boundary of the same project: instead of starting with explicit sections, the analyzer tries
to infer local sections from existing code and to transport them upward into the stack.

Theory packs are no longer a plugin mechanism with independent semantics. They are local
geometries inside the same project. Each pack says which sites matter, which solver or proof
technology can act locally there, and what sort of trust upgrade is justified when the local
problem is discharged.

The generation agent is no longer a free-floating LLM wrapper. It is a cover-growing and
cover-repairing machine. Its task is to produce local sections whose interfaces, proof
obligations, and semantics glue. Its failures are therefore not opaque generation mishaps;
they are obstructions in the stack.

\section*{Axioms of the project}
A project feels single when its commitments can be stated axiomatically. The repository is
already near that point. A reasonably faithful list of axioms is the following.
\begin{enumerate}[leftmargin=2em, itemsep=0.4em]
  \item \textbf{Locality first.} Software facts are obtained locally before they are claimed
  globally.
  \item \textbf{Types before heuristics.} Every reusable judgment should be represented, when
  possible, as a typed local section rather than a nameless heuristic.
  \item \textbf{Evidence should be explicit.} Trust levels, provenance, and residuals belong
  to the object, not merely to a log file.
  \item \textbf{Non-gluing is structured.} Bugs, contradictions, hallucinations, and failed
  synthesis attempts should be classified as organized obstructions.
  \item \textbf{Repair is refinement.} Failed globality should trigger stronger local models,
  finer covers, or better transports rather than silent fallback.
  \item \textbf{Adoption must be gradual.} The stack should admit both zero-change analysis and
  deeply annotated mixed-mode programming.
  \item \textbf{Generation belongs inside semantics.} Planning and synthesis are not beyond the
  theory; they are part of it.
\end{enumerate}

This axiomatic register is one of the main things the monograph was previously missing. Once
it is added, the later chapters stop sounding like catalog entries and start sounding like
consequences.

\section*{The repository as one mathematical object}
A final way to say the same thing is this: the repository is trying to build one mathematical
object with many interfaces. One interface is foundational and symbolic. One is operational
and proof-oriented. One is diagnostic and user-facing. One is heuristic and exploratory.
What gives the project dignity is not that each interface is already perfect. It is that each
interface can be described in terms of the same stack.

If the rest of the manuscript succeeds, the reader should reach the end feeling not that they
have toured many ambitious ideas, but that they have spent time inside one ambitious idea
that keeps differentiating itself into the forms demanded by software practice.
